\documentclass[10pt]{article}

\usepackage[margin=1in]{geometry}
\usepackage{graphicx}
\usepackage{booktabs}
\usepackage{hyperref}

\title{How Much Adaptation Is Enough?\\Lightweight Drift Correction for Brain-to-Text BCI}
\author{AdaptLadder-BCI}
\date{}

\begin{document}
\maketitle

\begin{abstract}
Brain-to-text BCI decoders are vulnerable to cross-session changes in neural feature
distributions. This short paper studies lightweight input adaptation as a first-line
response to session drift. We quantify drift across T15 copy-task sessions, test simple
normalization and moment-matching corrections, and identify when stronger adaptation
may be necessary.
\end{abstract}

\section{Introduction}
Cross-session drift can degrade BCI decoder reliability and increase recalibration
burden. Rather than immediately changing decoder weights, we ask how far simple
input-space corrections can go.

\section{Related Work}
To be filled with brain-to-text BCI, neural decoder adaptation, test-time adaptation,
and recalibration references.

\section{Data}
Primary analysis uses T15 copy-task neural features. The first sanity pass reports
session counts, trial counts, trial lengths, feature dimensionality, and available
labels/splits.

\section{Methods}
\subsection{Drift metrics}
For each session we estimate feature-wise mean, standard deviation, and diagonal
covariance. We report distance from the first session and from the previous session.

\subsection{Adaptation ladder}
We compare no adaptation, target-session z-scoring, and source-to-target moment
matching. Decoder-based metrics are planned as the main performance endpoint.

\section{Results}
Initial results will include T15 dataset summary, drift-over-time, PCA session maps,
and adaptation before/after metrics.

\section{Discussion}
The central question is whether lightweight correction is enough, and when the
remaining error calls for stronger adaptation.

\section{Limitations}
Distribution alignment alone is not a sufficient performance metric. WER, PER, CTC
loss, decoder confidence, or related decoder-facing metrics are required for a stronger
paper.

\section{Conclusion}
Adaptation ladders can provide a practical way to decide how much correction a
brain-to-text BCI session needs before moving to heavier recalibration.

\end{document}
